\begin{abstract}
    %Integrating recently obtained \ac{SAR} images and their processing consists a significant challenge in the domain of phase estimation, primarily due to the unprecedented volume  of \acs{SAR} images acquired with Sentinel-1  mission ($6-12$ days). This paper propose a novel approach to sequentially integrate each newly acquired image. The methodology is based on the \ac{PL}, or more precisely on the \ac{MLE}. The core concept involves dividing the data set into two blocks, one representing the previous images and the other representing the new additions. Leveraging past images as a prior, our approach sequentially estimate the phase of the newly acquired data. A crucial factor for the success of \acs{PL} algorithms lies in the estimation of the covariance matrix, and many \acs{PL} algorithms in the state of the art are mainly built on a Gaussian distribution assumption, however it tends to restrict the accuracy of covariance matrix estimation in high-dimensional \acs{SAR} images.  To enhance robustness against non-Gaussian data, we explore a scaled mixture of Gaussian model. The performances of the new approach is validated on simulations and Sentinel-1 \acs{SAR} images.
    %\ac{InSAR} techniques are widely used to monitore surface deformation that occurred between two dates, which is done through phase estimation. Treating each slave image according to the image master is not sufficient to achieve reliable results, that's why with a time series of \ac{SAR} images, an improvement was observed. \ac{MT-InSAR} exploits the covariance matrix of the \acs{SAR} images time series, and estimate the coherence matrix and the phases. However, these approaches become computationally expensive as the time series lenghtens and new components are added to the data vector. This paper propose a novel approach to sequentially integrate each newly acquired image. The methodology is based on the \ac{PL}, or more precisely on the \ac{MLE}. The core concept involves dividing the data set into two blocks, one representing the previous images and the other representing the new additions. Leveraging past images as a prior, our approach sequentially estimate the phase of the newly acquired data. A crucial factor for the success of \acs{PL} algorithms lies in the estimation of the covariance matrix, and many \acs{PL} algorithms in the state of the art are mainly built on a Gaussian distribution assumption, however it tends to restrict the accuracy of covariance matrix estimation in high-dimensional \acs{SAR} images. In urban areas, the data often fits a non-Gaussian distribution. To enhance robustness against non-Gaussian data, we explore a scaled mixture of Gaussian model. The performances of the new approach is validated on simulations and Sentinel-1 \acs{SAR} images.
    
    %\ac{InSAR} techniques monitor surface deformation between two dates via phase estimation. Treating each slave image with respect to the master image is not enough for reliable results, which is why improvements are seen with a time series of \ac{SAR} images. 
    \ac{MT-InSAR} exploits \ac{SARTS} for surface deformation monitoring via phase difference (with respect to a reference image) estimation. Most of the actual state-of-the-art \acs{MT-InSAR} rely on temporal covariance matrix of the \acs{SARTS}, assuming Gaussian distribution.
    %of the time series to estimate the coherence matrix and phases. 
    However, these approaches become computationally expensive when the time series lengthens and new images are added to the data vector. 
    This paper proposes a novel approach to sequentially integrate each newly acquired image using \ac{PL} and %\textcolor{blue}{
    \ac{MLE}.
    %}. 
    The methodology divides the data into blocks, using previous images and estimations as a prior to sequentially estimate the phase of the new image. Actually, this framework allows to consider non Gaussian distributions, such as a mixture of scaled Gaussian distribution, which is particularly important to consider when dealing with urban areas.  
    %A crucial factor for the success of \ac{PL} algorithms is the estimation of the covariance matrix, often based on a Gaussian distribution assumption. However, in urban areas, data often fits a non-Gaussian distribution. To enhance robustness, we explore a scaled mixture of Gaussian model. Compared to offline approaches, the proposed approach provides identical performance, but present significant computational efficiency.   %The performance of the new approach is validated on simulations and Sentinel-1 \ac{SAR} images.
    
    %\textcolor{red}{Multi-temporal Interferometry Synthetic Aperture Radar exploits \acs{SARTS} for surface deformation monitoring via phase difference (with respect to a reference image) estimation. Most actual state-of-the-art MT-InSAR approaches, e.g. phase linking, rely on covariance matrix estimation of \acs{SARTS}, assuming gaussian distribution for SAR data. These approaches become computationally expensive when the time series lengthens and new images are added to the data vector. 
    %This paper proposes a novel approach to sequentially integrate each newly acquired image using \ac{PL} and maximum likelihood estimator. In this approach, we model SAR data with a scaled mixture of Gaussian models in order to better represent the reality. The phase differences of new images are estimated sequentially based on prior information issued from previous images. The performance of the proposed approach is validated on both simulations and Sentinel-1 images. Compared to offline approaches, the proposed approach provides identical performance, but present significant computational efficiency.  }
\end{abstract}

\begin{IEEEkeywords}
    %Enter key words or phrases in alphabetical order, separated by commas. For a list of suggested keywords, send a blank e-mail to keywords@ieee.org or visit \url{http://www.ieee.org/organizations/pubs/ani_prod/keywrd98.txt}
    %\textcolor{blue}{\acl{MT-InSAR}, \acl{MLE}},
    \acl{MT-InSAR}, \acl{MLE}, covariance matrix estimation, mixture of Scaled Gaussian distribution
\end{IEEEkeywords}
%\blfootnote{The source code is available on GitHub at the following address: https://github.com/DanaElhajjar/S-SG-MLE-PL}